% main.tex — Critical Assessment of MULTING/IDM Framework
% Sergey Boyko · Ronin Institute · 2026-06-19
% NOT_VALIDATION · NOT_REFUTATION · OUR_RECONSTRUCTION
% Status: DRAFT · NOT_FOR_SUBMISSION · PARTLY SUPERSEDED (see below)
%
% =====================================================================
% SUPERSEDED IN PART BY FINDING B0 (2026-08-03) — READ BEFORE CITING
% ---------------------------------------------------------------------
% Last substantive edit of this file: 2026-07-22.
% Finding B0 postdates it by 12 days and is NOT yet incorporated
% throughout:
%   experiments/20260803-bridge/FINDING_table_a1_provenance.md
%
% B0, in the finding's own (deliberately hedged) wording: the material
% available to us does not contain a deterministic, published,
% reproducible operator F -> H that produced Table A1. Every column of
% Table A1 — including H-MULT and H-FLRW — was requested from an online
% AI service, per the author's own caption and surrounding text
% (preprint pp. 38-39, [VERIFIED-SOURCE]). The author states plainly
% that those responses are "not necessarily trustworthy".
%
% CONSEQUENCE FOR THIS FILE: every quantity here derived FROM Table A1
% — epsilon(z), q(z), w_eff, and the f-sigma8 shape correlation
% (Sec. \ref{ssec:fsig8}) — characterises that AI response, NOT MULTING.
% The results are not wrong; they are mis-subjected. Sections carrying
% affected quantities are flagged inline with \tableAsubject.
%
% Also superseded: data/table_a1_reported.csv (4 wrong fields in row 1,
% notably H_MULT 71.1 -> 73.0). Use data/table_a1_source_verified.csv.
%
% UNAFFECTED: Eq.32, 7:9:17, fermion masses, N_opt/chi2 (PDG/Planck
% based), and the kSZ track (ACT+SDSS, never touches Table A1).
% =====================================================================

\documentclass[twocolumn,prd,amsmath,amssymb,floatfix]{revtex4-2}

\usepackage{graphicx}
\usepackage{hyperref}
\usepackage{booktabs}
\usepackage{amsmath}
\usepackage{xcolor}
\usepackage{siunitx}
\usepackage{orcidlink}

% Safety label macro
\newcommand{\reconstruction}{\textcolor{gray}{\small[OUR\_RECONSTRUCTION]}}
\newcommand{\notvalidation}{\textcolor{gray}{\small[NOT\_VALIDATION]}}
% Subject-of-claim flag, added 2026-08-04 after finding B0: marks any
% quantity derived from Table A1, whose subject is the AI service's
% response rather than MULTING itself.
\newcommand{\tableAsubject}{\textcolor{gray}{\small[SUBJECT:
  TABLE~A1 AI RESPONSE, NOT MULTING --- see B0, 2026-08-03]}}

\begin{document}

\title{Systematic Assessment of the MULTING/IDM Gravitational Framework:\\
Confirmed Assets, Identified Gaps, and Independent Reconstructions}

\author{Sergey Boyko\,\orcidlink{0009-0009-2178-5701}}
\affiliation{Ronin Institute for Independent Scholarship}
\email{sergeikuch80@gmail.com}
\date{\today}

\begin{abstract}
We present a systematic 36-point assessment of the MULTING/IDM preprint
(Buckholtz 2026, \texttt{preprints202511.0598.v6}), which proposes a
multipole-expansion extension of Newtonian gravity (MULTING) and an
isomer-based dark-matter framework (IDM). Using PDG 2024 constants, Planck 2018
cosmological parameters, and the MCXC/XMM galaxy cluster catalogs, we
independently verify the mass relations in Equations~(31) and~(32), quantify
the statistical preference for MULTING vs.\ $\Lambda$CDM in H(z) data, and
derive formal constraints on the IDM isomer count from the baryon density ratio
$\omega_\mathrm{DM}/\omega_b$.

We find: (1)~the empirical boson-mass relation [7:9:17] holds to $<0.05\%$
in mass units (the scale-free ratio $m_W/m_H = \sqrt{7/17}$ holds to $0.04\%$,
though per-boson the $W$ deviates by $3.8\sigma$ and $H$ by $1.1\sigma$ when
$Z$-anchored); stated parameter-free, the relation is consistent only with the
heavier (CDF-II/CMS) boson masses ($\chi^2 = 2.0$, $p = 0.36$) and excludes the
light/Standard-Model set ($\chi^2 = 39$), making it a falsifiable discriminant
on the $W$-mass anomaly. The fermion-gravity link (Eq.~32) holds at $1.0\sigma$,
with a trials-corrected $p_\mathrm{look\text{-}elsewhere} = 0.16$ over an
explicitly declared expression space
--- both empirical assets independent of MULTING/IDM cosmology; (2)~the fermion mass spectrum
(Eqs.~21--24) reproduces the muon mass to $0.47\%$ and quark geometric means
to $<0.31\%$ using PDG 2024 constants; (3)~the MULTING dipole term is
rendered negligible at cluster scales for the tabulated $\beta_d = 4.5$,
requiring $\beta_d \gtrsim 10^{3.2}$ for a 1\% observational signature ---
formally quantified via Birge Ratio $R_B = 15.9$ ($\beta_d$) and $R_B = 24.1$
($\beta_q$), both with $p < 10^{-4}$;
(4)~equal-mass $N=5$ IDM isomers are excluded at $5.79\sigma$ by Planck 2018,
but five isomers with $\bar{m} = 1.074\,m_p$ are Planck-compatible (EXP-N);
(5)~the $\Delta N_\mathrm{eff}$ exclusion (Planck bound exceeded by a factor
of $\sim$80--215) applies to a thermalized mirror sector only;
gravitational-only IDM gives $\Delta N_\mathrm{eff}
\lesssim 10^{-40}$, trivially within Planck bounds, though this eliminates
the dark baryogenesis mechanism (EXP-O);
(6)~the $H_\mathrm{FLRW}$ baseline used in Table~A1 is inconsistent with standard
$\Lambda$CDM (Planck MAE $= 120$~km\,s$^{-1}$\,Mpc$^{-1}$); a power-law
$H(z) = 54.07\,(1+z)^{0.884}$ fits the 11 Table~A1 points with
MAE $= 5.04$~km\,s$^{-1}$\,Mpc$^{-1}$.

These results neither validate nor refute MULTING/IDM, but map precisely which
components rest on verified foundations versus which require additional physics.
\end{abstract}

\maketitle

% ──────────────────────────────────────────────
\section{Introduction}
\label{sec:intro}
% ──────────────────────────────────────────────

The framework of Buckholtz~\cite{buckholtz2026} proposes a multipole-expansion
extension of Newtonian gravity (MULTING) in which galaxy clusters exert a
dipole force component $F_d$ and a quadrupole component $F_q$ that supplement
the monopole $F_m$, together modifying the Hubble expansion rate.
The preprint provides a table of values (Table~A1) for
$H_\mathrm{MULT}(z)$ at 11 redshift points.
We initially treated this table as a MULTING calculation; that assumption
was audited and withdrawn on 2026-08-03 (finding~B0).
The author states in the caption and surrounding text that Table~A1 shows
``responses, to our prompt, by one online service that has bases in
artificial intelligence'', that the $H_\mathrm{MULT}$ column specifically
was obtained because ``the prompt asked for a value that MULTING
suggests'', and that the responses are ``not necessarily trustworthy''.
The material available to us therefore does not contain a deterministic,
published, reproducible operator carrying the force law into $H(z)$.
The author labelled the table accurately; the error was ours, in reading
it as a model calculation.
Results in this paper that were derived from Table~A1 are flagged
accordingly and characterise that response rather than the theory.
\tableAsubject

The six fundamental constants identified by Rees~\cite{rees1999} include
the ratio of electromagnetic to gravitational coupling strengths
$N \approx 10^{36}$ and the cosmological constant $\lambda \approx 0.7$.
MULTING/IDM addresses both: Eq.~(32) provides an empirical link between
$N$ and lepton masses ($1.0\sigma$, \S\ref{ssec:eq32}), while the MULTING
dipole term offers a mechanical alternative to $\lambda$ at first
cosmological perturbation order (\S\ref{sec:hz}).

A separate and timely test arises from the boson-mass relation Eq.~(31). Because
it fixes the ratios $m_W/m_Z$ and $m_H/m_Z$ with no free parameter, it speaks
directly to the unresolved $W$-mass tension --- the CDF-II~2022 value lies
$\sim 7\sigma$ above the Standard-Model expectation, while recent LHC
measurements favor the lighter, SM-consistent value. As we show
(\S\ref{ssec:bosons}), Eq.~(31) is statistically consistent only with the
heavier measurement set, turning a numerological-looking ratio into a concrete,
falsifiable prediction that forthcoming $W$-mass combinations will confirm or
exclude.

Beyond the mass relations, the preprint proposes that the thermal energy of
intergalactic medium at cosmic-web nodes constitutes a second, scalar gravitational
source that supplements the cluster monopole.
This mechanism occupies a gap that competing approaches to anomalous gravity
leave unfilled: Verlinde's entropic formulation~\cite{verlinde2016} reproduces
galaxy-scale rotation curves but is incompatible with cluster-scale lensing
observations~\cite{hossenfelder2017,ettori2016} and with the Bullet
Cluster~\cite{clowe2006}.
As a preliminary empirical test we searched for the predicted correlation between
the lensing-to-hydrostatic mass gap ($\Delta M = M_\mathrm{WL} - M_\mathrm{HE}$)
and an ICM thermal energy proxy ($M_\mathrm{gas}\times T_x$) in the CCCP
sample~\cite{mahdavi2013} ($N=50$, $0.15<z<0.55$).
The Pearson correlation is $r=0.021$ ($p=0.883$) --- the sub-hypothesis that
\emph{cluster interior} thermal energy predicts the mass gap is not supported.
A secondary, unanticipated pattern emerged from the same sample: conditioning on
$M_\mathrm{WL}$, the partial correlation between $\Delta M$ and the thermal
energy proxy is $r=-0.701$ ($p<10^{-4}$) --- at fixed lensing mass, clusters with
\emph{more} ICM thermal energy show a \emph{smaller} mass gap. We tested three
standard-physics explanations for this partial correlation and none mediates it:
splitting the sample at the median $M_\mathrm{WL}$ leaves the
partial correlation unchanged in both subsamples ($r=-0.732$ vs.\ $r=-0.731$,
Fisher $z=0.005$, $p=0.50$, ruling out a mass-threshold effect); controlling
additionally for X-ray centroid shift $w_X$ (a standard cluster
morphology/relaxation indicator, available for 47/50 clusters from the same
catalog) leaves the partial correlation essentially unchanged
($r=-0.726$ vs.\ $-0.714$ baseline on the same subsample), ruling out
dynamical-state confounding; and controlling for central entropy $K_0$ (a
standard AGN radio/kinetic-feedback and cool-core proxy) similarly leaves it
unchanged ($r=-0.718$ vs.\ $-0.714$ baseline), ruling out AGN-feedback
confounding. However, none of these three tests controls for X-ray temperature
$T_X$ itself, and $T_X$ enters the correlation on \emph{both} sides by
construction: directly in the thermal-energy proxy ($M_\mathrm{gas}\times T_x$),
and indirectly in $\Delta M$ through $M_\mathrm{HE}$, since standard
hydrostatic-equilibrium mass estimation requires the temperature profile as an
input. A dedicated test controlling for both $M_\mathrm{WL}$ and $T_x$
simultaneously finds that $M_\mathrm{gas}$ alone --- the thermal-mass
component not multiplicatively tied to $T_x$ --- no longer significantly
correlates with $\Delta M$ ($r=-0.08$, $p=0.58$, bootstrap 95\% CI
$[-0.39,0.27]$, $N=50$), while $T_x$ alone at fixed $M_\mathrm{WL}$ is the
single strongest predictor found ($r=-0.81$). We therefore do
\emph{not} characterize this partial correlation as an unexplained empirical
pattern robust to standard physics, but we also do not identify its mechanism:
two explanations remain equally consistent with the data and are not
distinguished by this test --- a shared, largely definitional dependence on
cluster temperature (via $M_\mathrm{HE}$'s construction), or a genuine common
physical driver such as cluster dynamical state affecting both $T_x$ and the
hydrostatic-equilibrium assumption independently. Resolving this would require
either an independent, temperature-decoupled mass estimate or a
dynamical-state stratification of the existing sample; neither has been done.
It does not constitute evidence for the preprint's mechanism under any reading
considered here: it is a cluster-interior (ICM) result, not the cosmic-web
filament (WHIM) effect the mechanism proposes, and no causal model for the
partial correlation itself is offered here. Crucially, the preprint's
mechanism refers to warm-hot
intergalactic medium (WHIM, $T\sim 10^5$--$10^7$~K) in cosmic-web
\emph{filaments} rather than the cluster ICM --- a physically distinct gas
phase and radius range not subject to this specific temperature-sharing
concern; the broader hypothesis remains inconclusive pending WHIM-specific
observations.
All hypothesis tests in this paper are labelled \reconstruction\ and follow
the Falsification Ladder protocol of our supplement.

The principal challenge for external assessment is that the bridge connecting
cluster physical parameters ($k_A$, $r_A$, $D_{C:AB}$) to $H_\mathrm{MULT}$
is not derivable from first principles using standard cluster physics: the
tabulated $\varepsilon(z) \equiv (H_\mathrm{MULT}/H_\mathrm{FLRW})^2 - 1$
peaks at $z = 0.40$ and exhibits a secondary rise at $z = 8.50$, a
non-monotone shape inconsistent with virial-theorem cluster energetics (NR-004;
NR-008).  The phenomenological parameters $\beta_d$ and $\beta_q$ vary by
factors of $5.8\times$ and $95\times$ across three AI services queried in the
preprint, indicating that these parameters are not uniquely determined by the
stated procedure.

Our contribution is to separate what is independently verifiable from what
requires the author's undisclosed numerical procedure.
Section~\ref{sec:assets} verifies the mass relations (Eqs.~31, 32) against
PDG 2024.
Section~\ref{sec:hz} quantifies the statistical comparison of MULTING and
$\Lambda$CDM on H(z) data and identifies the $\beta$-rescaling gap.
Section~\ref{sec:idm} audits the IDM isomer-count prediction against Planck
2018 and standard dark-matter constraints.
All numerical results carry the label \reconstruction\ to distinguish our
independent reconstructions from author-attributed claims.

% ──────────────────────────────────────────────
\section{Verified Assets}
\label{sec:assets}
% ──────────────────────────────────────────────

\subsection{Boson Mass Relation (Eq.~31)}
\label{ssec:bosons}

The relation $m_W^2 : m_Z^2 : m_H^2 = 7:9:17$ (Eq.~31 of~\cite{buckholtz2026})
was tested against PDG 2024 masses~\cite{pdg2024}:
$m_W = 80.369 \pm 0.013$~GeV (PDG 2024 average, CDF-II 2022 excluded;
its inclusion drops the combination compatibility to $0.5\%$),
$m_Z = 91.1876 \pm 0.0021$~GeV,
$m_H = 125.20 \pm 0.11$~GeV (PDG 2024).

Using the $Z$ mass as anchor ($\lambda = m_Z^2/9 = 923.9$~GeV$^2$), the
predicted Higgs mass is $125.33$~GeV, agreeing with the measured value at
$1.14\sigma$.  The predicted $W$ mass is $80.420$~GeV, deviating by
$3.81\sigma$ from the PDG 2024 average.  Both predictions lie at the heavier
end of the currently discrepant measurements: the $W$ prediction agrees with
the CDF-II~2022 value ($80.4335 \pm 0.0094$) at $1.45\sigma$, and the $H$
prediction agrees with the CMS~2022 value ($125.35 \pm 0.15$) at $0.16\sigma$.
We emphasize that the per-boson significance is anchor-dependent: fixing the
scale $\lambda$ to any one boson renders that boson exact by construction
($0\sigma$), so the honest figure of merit is the \emph{largest} residual
across the three masses --- here the $W$ boson at $\approx 3.8\sigma$ --- not
the best-fitting one. We therefore report the relation as precise in mass
units ($<0.05\%$, below) but only marginally consistent in error-weighted
units, dominated by the $W$ tension.
\reconstruction

The anchor dependence is removed entirely by stating the relation in
\emph{scale-free} ratio form, which carries the same content without choosing
a reference boson:
\begin{equation}
    \frac{m_W}{m_H} = \sqrt{\tfrac{7}{17}} = 0.6417, \qquad
    \frac{m_W}{m_Z} = \sqrt{\tfrac{7}{9}} = 0.8819 .
\end{equation}
The ratio $m_W/m_H$ holds to $0.04\%$ against the PDG 2024 averages, and to
$0.003\%$ against the heavier $(m_W^\mathrm{CDF\text{-}II}, m_H^\mathrm{CMS})$
pair --- consistent with the observation above that the relation favors the
heavier end of both discrepant measurements. Because $m_W/m_H$ is independent
of the scale anchor, it is the cleaner observable: future $W$-mass combinations
provide a direct, anchor-free falsification test of Eq.~31.
\reconstruction

\paragraph{Preference statistic.}
Stated scale-free, Eq.~31 fixes the two ratios $m_W/m_Z=\sqrt{7/9}$ and
$m_H/m_Z=\sqrt{17/9}$ with \emph{zero} free parameters. We compute the 2-dof
$\chi^2$ of these predictions against three measurement scenarios that differ
only by the currently-discrepant $W$ and $H$ values
(Table~\ref{tab:c6pref}). The relation is statistically consistent
($\chi^2=2.0$, $p=0.36$) only with the \emph{heavy} set
(CDF-II~$m_W$ + CMS~$m_H$); the PDG-2024 average gives $\chi^2=15.6$
($p=4\times10^{-4}$) and the light set (CMS-24~$m_W$ + ATLAS~$m_H$) is excluded
at $\chi^2=39$. Eq.~31 therefore acts as a falsifiable discriminant favoring the
heavier, outlier measurements rather than the experimental consensus; it is a
bet, not evidence, and a future $W$ world-average converging below
$\sim 80.38$~GeV would exclude its error-weighted support at $>3\sigma$.
\reconstruction

\begin{table}[h]
\centering
\caption{7:9:17 zero-parameter $\chi^2$ (2 dof) by measurement scenario.}
\label{tab:c6pref}
\begin{tabular}{lccc}
\hline
Scenario & $\chi^2$ & $p$ & verdict \\
\hline
heavy (CDF-II $W$ + CMS $H$) & 2.0 & 0.36 & consistent \\
PDG-2024 average & 15.6 & $4\times10^{-4}$ & disfavored \\
light (CMS-24 $W$ + ATLAS $H$) & 39.0 & $<10^{-5}$ & excluded \\
\hline
\end{tabular}
\end{table}

A least-squares scale fit over all three masses simultaneously gives
$\lambda_\mathrm{LSQ} = 923.0$~GeV$^2$, with percentage deviations of
$-0.006\%$ ($W$), $+0.048\%$ ($Z$), $-0.012\%$ ($H$) in mass units ---
all below $0.05\%$, which is remarkably precise for a simple integer ratio
involving three independently measured bosons.
The ladder is specific to the electroweak bosons: the top quark gives
$(m_t/(m_Z/3))^2 = 32.3$, off the nearest integer ($32$) by $0.87\%$
($2.5\sigma$) --- an order of magnitude worse than the $W/Z/H$ fits --- so
Eq.~31 does not extend to the heavy-fermion sector.

\paragraph{Look-elsewhere for integer assignment.}
A systematic scan over all integer pairs $(n_W, n_H)$ with $1 \le n_W, n_H \le 50$
finds that $(7, 17)$ is the \emph{unique} pair that brings both $m_W^2$ and $m_H^2$
within any threshold from $0.2\%$ up to $2\%$ of integer multiples of $(m_Z/3)^2$.
No alternative integer pair exists in this range at any tested threshold.
A Monte Carlo over $5 \times 10^5$ draws of hypothetical boson masses uniformly
distributed in $\pm5\%$ windows around the measured values shows that the probability
of an arbitrary mass pair landing within $0.2\%$ of \emph{any} integer assignment
is approximately $1/600$.

A further anchor-sensitivity scan over units $(m_Z/k)^2$ for $k = 2, \ldots, 10$
confirms that $k=3$ is the \emph{only} anchor value among all tested denominators
for which both $m_W^2$ and $m_H^2$ simultaneously fall within $2\%$ of integers
$(n_W, n_H)$ with $n \le 50$.
For $k \ne 3$, at most one of the two masses is close to an integer in these units.
This anchor-uniqueness result addresses the concern that the integer pattern could be
reproduced with alternative normalization choices.\reconstruction

\subsection{Fermion-Gravity Link (Eq.~32)}
\label{ssec:eq32}

The empirical relation
\begin{equation}
    \frac{4}{3}\left(\frac{m_\tau}{m_e}\right)^{12} = \frac{\alpha_\mathrm{EM}}{\alpha_G}
    \label{eq:c9}
\end{equation}
where $\alpha_G = G m_e^2 / (\hbar c)$, was evaluated using PDG 2024
constants~\cite{pdg2024}: $m_\tau = 1776.93 \pm 0.09$~MeV, $m_e = 0.51100$~MeV,
$\alpha_\mathrm{EM}^{-1} = 137.036$, $G = 6.674 \times 10^{-11}$~m$^3$~kg$^{-1}$~s$^{-2}$.
(Corrected 2026-07-11: an earlier version of this subsection used
$m_\tau = 1776.86 \pm 0.12$~MeV mislabeled ``PDG 2024''; that is the
PDG-2022-era value. The relation survives the correction, now at
$1.0\sigma$ rather than $0.17\sigma$.)

We find LHS/RHS~$= 1.000608$, corresponding to a $1.0\sigma$ agreement
dominated by the $m_\tau$ uncertainty.
\reconstruction

\paragraph{Trials factor (added 2026-08-10).}
The $1.0\sigma$ figure quantifies the experimental uncertainty \emph{given} the
formula.  It is not the probability of finding such a relation somewhere in an
admissible expression space, and that second quantity had not been computed.
We therefore applied to Eq.~(\ref{eq:c9}) the same enumeration used to reject a
proposed geometric origin for its $4/3$ prefactor (NR-009).
The space, declared before scanning, is all
$p/q\,(m_A/m_B)^n$ with $p,q \le 12$ in lowest terms, $m_A/m_B > 1$ drawn from
$\{e,\mu,\tau,u,d,s,c,b,t,W,Z,H,p\}$, and integer $n \le 24$: 170\,352
expressions.
Of these, 35\,672 could be \emph{stated} to better than $0.5\%$ once the PDG
uncertainties of their own mass inputs are propagated --- the filter matters,
because the nearest unfiltered competitor,
$12\,(m_c/m_u)^{15}$ at $0.057\%$, carries a $265\%$ uncertainty of its own
(light-quark masses raised to the fifteenth power) and could never be claimed as
a precision relation.
Within that statable space Eq.~(\ref{eq:c9}) is the \emph{only} expression
reaching its own $0.059\%$.
Drawing 4000 random targets log-uniformly within a decade of
$\alpha_\mathrm{EM}/\alpha_G$ and recording the best hit available to each from
the same space gives a median best hit of $0.228\%$ and a 10th percentile of
$0.037\%$; Eq.~(\ref{eq:c9})'s $0.059\%$ corresponds to
\begin{equation}
  p_\mathrm{look\text{-}elsewhere} = 0.16 ,
\end{equation}
i.e.\ roughly one arbitrary target in six does at least as well.
The relation is twice better than the median and is not thereby refuted; its
agreement is, however, \emph{ordinary} rather than extraordinary for a target of
this magnitude, and the Belle~II test below remains falsifiable independently of
how the formula was found.
The $p$-value is conditional on the declared grammar: a wider one raises the
trials factor and lowers it.
\reconstruction

\paragraph{Formal uncertainty propagation.}
A first-order propagation through Eq.~(\ref{eq:c9}) using PDG 2024 and
CODATA 2018 uncertainties gives:
\begin{align}
  \sigma_\mathrm{LHS}(m_\tau) &= \tfrac{4}{3}\cdot 12\cdot
    \left(\frac{m_\tau}{m_e}\right)^{11}\!\frac{\sigma_{m_\tau}}{m_e}
    = 0.061\%\text{ of LHS} \notag \\
  \sigma_\mathrm{RHS}(G) &= \frac{\sigma_G}{G}\,\mathrm{RHS}
    = 0.0022\%\text{ of RHS} \notag \\
  \sigma_\mathrm{combined} &= 0.061\%\text{ of RHS};\quad
    |\mathrm{LHS}-\mathrm{RHS}|/\sigma = 1.00\sigma .
\end{align}
The dominant uncertainty is \emph{not} $G$ but $m_\tau$: the exponent 12
amplifies the $m_\tau$ contribution by a factor of $12^2 = 144$ relative to $G$'s
contribution at the same relative precision.
With the Belle~II target $\sigma_{m_\tau} = 0.01$~MeV~\cite{belleii2019},
the combined significance would rise to $8.54\sigma$ — at the current
central value, a Belle~II-precision measurement would decisively
\emph{disfavor} exact equality unless the tau mass itself shifts toward
$1776.840$~MeV (Sec.~\ref{ssec:eq32}'s companion note gives the two-sided
reading of this test in full).
Partial derivatives verified against finite differences to $10^{-7}$ precision.
\reconstruction

\paragraph{Belle~II has already measured $m_\tau$ (added 2026-09-03).}
The projection above treated Belle~II as a purely future test; it is not.
Belle~II's own 2023 measurement~\cite{belleii2023taumass} gives
$m_\tau = 1777.09 \pm 0.08\,\mathrm{(stat)} \pm 0.11\,\mathrm{(syst)}$~MeV
(combined $0.136$~MeV). Checked directly against Eq.~(\ref{eq:c9})'s exact
prediction of $1776.840$~MeV, this single, genuinely independent
measurement — unlike the PDG-2024 world average used above, which itself
already folds in this and other measurements, and is therefore not an
independent check of an anchor built against an earlier world average —
gives
\begin{equation}
  |m_\tau^\mathrm{Belle~II} - m_\tau^\mathrm{Eq.~(\ref{eq:c9})}| / \sigma
  = 1.84\sigma .
\end{equation}
This is neither a clean confirmation nor a decisive rejection: it sits
between the $1.0\sigma$ PDG-world-average figure above and the
$8.54\sigma$ hypothetical target-precision projection, and should be
read as the honest current state of the direct experimental test, not
as a resolved result. Belle~II Run~2 (data-taking since 2024-02-20) may
sharpen this tension in either direction.
\reconstruction

\paragraph{Sensitivity to $G$ measurements.}
Different precision measurements of $G$ yield deviations of Eq.~(\ref{eq:c9})
ranging from $+0.049\%$ \cite{luo2009} to $+0.079\%$ \cite{quinn2013},
now at $0.8$--$1.3\sigma$ significance (versus the combined uncertainty of
Sec.~\ref{ssec:eq32} above). The spread among published $G$ values
($479$~ppm peak-to-peak across 10 experiments since 1982) is now smaller
than the deviation of Eq.~(\ref{eq:c9}) from equality ($608$~ppm with
CODATA 2018~\cite{codata2018}): unlike in the earlier, PDG-2022-era
analysis, no published $G$ determination renders Eq.~(\ref{eq:c9}) exact to
within its own quoted precision, and the tau-mass uncertainty is now
unambiguously the limiting factor for this test, consistent with the
amplification-factor argument above. The value of $G$ that makes
Eq.~(\ref{eq:c9}) exact is $G^* = 6.6702 \times 10^{-11}$~m$^3$\,kg$^{-1}$\,s$^{-2}$,
about $27\,\sigma_G$ from the CODATA 2018 central value using $G$'s own
(22~ppm) precision alone --- unsurprising, since the actual gap between the
two sides is set by the far looser tau-mass uncertainty, not by $G$. A
measurement at MICROSCOPE-class precision ($\sim\!1$~ppm in $G$) would
still leave Eq.~(\ref{eq:c9}) limited by $m_\tau$, not $G$, because the
$m_\tau$ uncertainty contributes $0.061\%$ regardless of $G$ precision.
\reconstruction

\paragraph{Context and status.}
Eq.~(\ref{eq:c9}) belongs to a recognized tradition of empirical relations
among lepton masses and fundamental constants --- the Koide
formula~\cite{sumino2009koide,xing2006koide,lima2006koide} and the
large-number coincidences first noted by Dirac~\cite{dirac1937}, recently
reviewed in the context of varying constants by Uzan~\cite{uzan2025constants}
and connected to mass ratios and gravity by Singh~\cite{singh2022massratios}.
Eq.~(\ref{eq:c9}) is distinct from these: whereas the Koide relation is a
permutation-symmetric sum rule over $\{e,\mu,\tau\}$, Eq.~(\ref{eq:c9}) ties
the single ratio $m_\tau/m_e$ to the electromagnetic-to-gravitational coupling
ratio $\alpha_\mathrm{EM}/\alpha_G$. A direct numerical test of Koide
extensions --- scanning $K^p \times (m_\tau/m_e)^n$ for all integer $p \in [-3,3]$
and $n \in [6,20]$, where $K = (m_e+m_\mu+m_\tau)/(\sqrt{m_e}+\sqrt{m_\mu}+\sqrt{m_\tau})^2
\approx 2/3$ is the Koide ratio --- finds no combination within $10\%$ of
$\alpha_\mathrm{EM}/\alpha_G$; the nearest miss is $K^{-1}(m_\tau/m_e)^{12}$
at $12.5\%$ deviation, confirming that Eq.~(\ref{eq:c9}) is not a Koide extension.
A systematic survey of 77 papers in hep-ph, hep-th, and math-ph addressing
lepton mass hierarchies, coupling constant ratios, and exceptional Jordan algebras
did not locate this specific relation; Targeted searches of both indexed
open-access databases and the open web support the same conclusion,
though related lepton-mass relations (e.g.~Singh's octonionic
derivation~\cite{singh2022massratios}) are actively studied.
We stress that Eq.~(\ref{eq:c9}) is presented as an \emph{empirical regularity,
mechanism open}: unlike the Koide formula, which admits a family-gauge-symmetry
derivation~\cite{sumino2009koide}, no first-principles derivation of the
prefactor $4/3$ or the exponent $12$ is available here. In particular, we
tested and could not support a geometric ($S^3$ curvature-balance) origin for
the $4/3$ factor: writing $4/3 = (n+1)/n$ and $12 = n(n+1)$ at $n=3$ is one of
several hundred equally simple two-function relabelings consistent at $n=3$,
so the geometric reading carries no predictive content beyond the numerical
coincidence itself.
\reconstruction

\subsection{Fermion Mass Spectrum (Eqs.~21--24)}
\label{ssec:fermions}

The IDM mass formula predicts lepton masses as
$m_k/m_e = (m_\tau/m_e)^{(k + \sigma_k \delta)/3}$
with $k \in \{0,2,3\}$ for $\{e, \mu, \tau\}$, $\sigma_k \in \{0,-1,0\}$,
and $\delta = 0.03668$ (Eqs.~21--24 of~\cite{buckholtz2026}).
Using PDG 2024 masses ($m_e = 0.51100$, $m_\tau = 1776.93$~MeV), the electron
is reproduced exactly by construction; the muon prediction is $106.16$~MeV
vs.\ PDG $105.66$~MeV (deviation $0.47\%$).
The best-fit $\delta = 0.03841$ reduces this to $< 0.01\%$ for the muon
while preserving the tau anchor.
\reconstruction

For quarks, the formula is applied to geometric means of up/down-type pair masses.
The three predictions agree with PDG 2024 geometric means to $< 0.31\%$
(lightest: $\sqrt{m_u m_d} = 3.21$ vs.\ predicted $3.20$~MeV;
intermediate: $\sqrt{m_s m_c} = 344.7$ vs.\ $345.2$~MeV;
heavy: $\sqrt{m_b m_t} = 850$ vs.\ $851$~MeV).

The inflaton candidate (Eqs.~27--30 with $M'=T'=1$) predicts
$m_\mathrm{inflaton} = m_Z/3 = 30.396$~GeV.
No candidate particle at this mass has been reported at LEP or LHC;
this constitutes an open experimental prediction.

% ──────────────────────────────────────────────
\section{MULTING H(z) Test}
\label{sec:hz}
% ──────────────────────────────────────────────

\subsection{Statistical Comparison with $\Lambda$CDM}

We compare MULTING against the 27-point Moresco~\cite{moresco2022} cosmic-chronometer
compilation using 5-fold cross-validation on the polynomial model
$H(z) = A(1+z)^2 + B(1+z)^3 + C(1+z)^4$, which mirrors the MULTING
force-law sign pattern $(+,-,+)$ for (monopole, dipole, quadrupole).
\reconstruction

Across all 5 folds the coefficient $B$ is consistently negative
(cross-validation coefficient of variation $\mathrm{CV}(B) = 3\%$),
confirming that the sign pattern is a stable data-driven feature and not a
fitting artefact.
The $\Delta\mathrm{AIC}$ between the MULTING polynomial and flat $\Lambda$CDM
is $+0.74$ (MULTING disfavoured), below the distinguishability threshold
$|\Delta\mathrm{AIC}| = 2$.
Out-of-sample performance on held-out CC points gives
$\chi^2_\mathrm{test} = 1.80$ (MULTING) vs.\ $1.71$ ($\Lambda$CDM):
statistically indistinguishable.
\reconstruction

This confirms the analytic result of Blanchet \& Le Tiec~\cite{blanchet2008,blanchet2009}:
a dipolar dark-matter fluid is degenerate with $\Lambda$CDM at first
cosmological perturbation order.
Distinguishing MULTING from $\Lambda$CDM on H(z) data alone therefore requires
either second-order observables or a first-principles derivation of $\beta_d$
that yields $\varepsilon \gtrsim 1\%$ (see \S\ref{ssec:beta}).

\subsection{The $\beta$-Rescaling Gap}
\label{ssec:beta}

The dipole force fraction at cluster scale is
\begin{equation}
    \varepsilon \equiv \frac{F_d}{F_m} = \beta_d \frac{k_A}{M} \frac{r_A}{D},
    \label{eq:epsilon}
\end{equation}
where $k_A = E_\mathrm{ICM}/c^2$ is the cluster kinetic mass and $D$ is the
inter-cluster distance (not the cluster radius), so $r_A/D \ll 1$ for any
physically meaningful pair. Evaluated on the real MCXC/PSZ2 catalogue
% [PROVENANCE RESOLVED 2026-08-10] The 2026-08-04 flag is lifted. The filter
% is recorded in data/pearson_r_test_results.md: clusters with E_thermal/c^2
% available (Path B, Y_SZ), i.e. data/clusters_clean.csv restricted to
% Ethermal_c2_Msun.notna(). That subset is exactly 548 of 1740 rows, and both
% medians reproduce from it: k_A/M = 1.739e-6 and R500/D = 9.979e-3 at
% D = 100 Mpc (median R500 = 0.998 Mpc). Verified by direct recomputation.
($n=548$, median $k_A/M = 1.7\times10^{-6}$, median
$r_A/D = 1.0\times10^{-2}$ at $D=D_0=100$~Mpc):

With the tabulated $\beta_d = 4.5$ (Table~A1 of~\cite{buckholtz2026}):
median $\varepsilon \approx 1.6 \times 10^{-7}$ (negligible). An observational
signature of $\varepsilon \gtrsim 1\%$ requires $\beta_d \gtrsim 2.9 \times 10^5$.
The AI-service run reported $\beta_d \approx 2\times10^4$ (Gemini), which yields
$\varepsilon \approx 0.07\%$ --- still fully negligible, not the ``marginally
non-negligible'' regime a naive typical-cluster estimate (using $r_A/D\sim1$,
i.e.\ conflating the cluster radius with the inter-cluster distance) would
suggest.
\reconstruction

To quantify the inter-service consistency of $\beta$ extractions, we apply the
Birge Ratio~\cite{birge1932} $R_B = \chi^2/(n-1)$ to three independent AI services
(Claude/NotebookLM: $\beta_d = 4.50$, $\beta_q = 18.0$;
Gemini: $4.25$, $8.10$; ChatGPT: $0.78$, $0.19$),
assuming a conservative relative uncertainty of $20\%$ per service.
We obtain $R_B = 15.9$ ($\chi^2 = 31.7$, $p < 10^{-4}$) for $\beta_d$ and
$R_B = 24.1$ ($\chi^2 = 48.3$, $p < 10^{-4}$) for $\beta_q$.
To achieve $R_B = 1$ (consistent services), a relative uncertainty of $80\%$
($\beta_d$) or $98\%$ ($\beta_q$) would be required --- far exceeding
the $\sim 20\%$ LLM extraction noise.
This \emph{Birge Ratio for AI Replication} (BRAI) confirms that $\beta$ values
are not uniquely determined by the stated procedure, consistent with the
author's own description of them as ``phenomenological computation''
whose derivation ``remains to be determined''.
\reconstruction

\subsection{Correlation with Real Cluster Catalogs: MULTING vs.\ Trivial Baselines}
\label{ssec:pearson}

The BRAI inconsistency quantified above is not merely a matter of differing
point estimates: the three AI-service $\beta$ pairs yield materially
different observational outcomes when the full MULTING force law is tested
against real cluster catalogs, independent of Table~A1.

\textbf{Method.} We compute
$\phi(z) = M_{500}/D^2 - 2\beta_d k_A r_{500}/D^3 + (\beta_q k_A r_{500})^2/D^4$,
$D \equiv D_0/(1+z)$ ($D_0 = 100$~Mpc; $D(z)$ is a phenomenological hypothesis,
not a derivation --- the same open bridge problem noted for $\beta_d,\beta_q$
in \S\ref{ssec:beta}), and $H_\mathrm{MULT}(z) = H_\mathrm{anchor}
\sqrt{\phi(z)/\phi(z_\mathrm{ref})}$, then correlate against 27
Moresco~\cite{moresco2022} cosmic-chronometer $H_\mathrm{CC}(z)$ points.
Cluster masses $M_{500}$ and radii $r_{500}$ are drawn from the
MCXC~\cite{mcxc2011} meta-catalogue; the kinetic-energy proxy $k_A = E_\mathrm{ICM}/c^2$
is estimated via PSZ2 $Y_\mathrm{SZ}$ under a $\Lambda$CDM angular-diameter
distance (Path~B).
\reconstruction

\textbf{Full-sample result.} On the flux-limited MCXC$+$PSZ2 sample ($n=443$
cluster--CC pairs after merger exclusion), the tabulated $\beta_d=4.5,
\beta_q=18.0$ gives $r=0.62$; a $\beta_d,\beta_q$ grid search
($\beta_d,\beta_q \in [10^2,10^8]$) finds no better fit than $r=0.6235$.
Both are beaten by a \emph{zero-parameter} monopole-only baseline
($\beta_d=\beta_q=0$, i.e.\ plain Newtonian $M_{500}/D^2$ scaling, no
MULTING-specific structure at all): $r = 0.7334$. This is in turn beaten by
a naive trivial model $H = H_0(1+z)^\alpha$ ($r=0.8904$) and by Planck-2018
$\Lambda$CDM~\cite{planck2018} ($r=0.8795$). The gap between monopole-only
and the $\beta_d=4.5,\beta_q=18.0$ fit is significant
(Fisher $z=3.05$, $p=0.0012$, one-tailed); AIC moderately favors $\Lambda$CDM
over the full MULTING formula ($\Delta\mathrm{AIC}=+2.5$ on the 27-point CC
compilation).
\reconstruction

\textbf{Mechanism.} Scanning $\beta_d$ and $\beta_q$ independently reveals
why: $r$ is essentially insensitive to $\beta_d$ across
$\beta_d \in [0, 10^6]$ at fixed $\beta_q$ (the dipole term is numerically
negligible for real cluster masses and radii at any tested magnitude), while
$r$ falls monotonically from $0.7334$ at $\beta_q=0$ to a hard floor of
$0.6234$--$0.6235$ by $\beta_q \sim 10$, remaining flat thereafter up to
$\beta_q=10^8$. Correlating each term's shape with $H_\mathrm{CC}(z)$ in
isolation gives a clean ranking: monopole alone $r=0.7334$ $>$ dipole-shape
alone $r=0.6655$ $>$ quadrupole-shape alone $r=0.6235$. The three AI-service
$\beta_q$ values from \S\ref{ssec:beta} sit on opposite sides of this
saturation threshold and land accordingly: ChatGPT ($\beta_q=0.19$, pre-saturation)
gives $r=0.7317$, near the zero-parameter optimum; Gemini ($\beta_q=8.10$)
and Claude/Table~A1 ($\beta_q=18.0$), both past threshold, give the same
degraded $r=0.6234$--$0.6235$ despite a $6\%$ difference in $\beta_d$ ---
consistent with $\beta_d$'s established inertness. This indicates the three
``AI thought-experiment'' fits did not converge toward a shared best fit;
their observational outcome is set almost entirely by which side of the
$\beta_q \sim 10$ saturation threshold each service happened to land on.
\reconstruction

\textbf{Cross-check on an unbiased subsample.} The MCXC test above is
flux-limited: a spurious $T_i$--$z$ correlation from Malmquist bias
($r(\log T_i, z) = 0.65$ on the full sample) is a priori indistinguishable
from a genuine MULTING signal. We repeat the test on the CHEX-MATE
combined catalogue (Planck-SZ--selected, closer to volume-limited), using
real XMM spectroscopic temperatures where available~\cite{rossetti2024} and
velocity-dispersion dynamical masses otherwise~\cite{sereno2025}. On the
$n=22$ subsample with real XMM $T_X$ (near-zero residual selection bias,
$r(\log M_{500}, z) = 0.094$), the trivial baseline still wins:
$r_\mathrm{trivial} = 0.9142$ vs.\ $r_\mathrm{MULT} = 0.7029$. On the full
$n=100$ CHEX-MATE sample the pattern repeats: $r_\mathrm{trivial}=0.9321$
vs.\ $r_\mathrm{MULT}=0.6269$. Because the unbiased subsample shows the same
gap as the flux-limited one, Malmquist bias is not the sole explanation for
MULTING underperforming the trivial baseline.
\reconstruction

\textbf{Interpretation.} The positive correlation ($r>0$ throughout)
confirms the correct qualitative sign of the MULTING prediction and is not,
by itself, contradicted. What is established is narrower: on every tested
sample (flux-limited MCXC, unbiased XMM, and full CHEX-MATE), adding the
dipole and quadrupole terms to the monopole never improves, and typically
degrades, the correlation with real $H(z)$ data relative to a
zero-parameter Newtonian baseline. Two explanations were considered: (a)
the quadrupole structure is genuinely disfavored by cluster data at these
scales, or (b) the $Y_\mathrm{SZ}$- and $\sigma_v$-derived $T_i$ proxies
used here are intrinsically noisier tracers of the cosmic expansion trend
than $M_{500}$ alone. A direct test partially distinguishes them: reshuffling
which cluster's $(k_A, r_{500})$ pair is attached to which $(z, M_{500})$
over 200 trials, holding the formula and $\beta_q=18$ fixed, shows the true
pairing beats every one of the 200 shuffled trials ($r=0.6234$ vs.\
shuffled mean $r^2=0.032$, maximum shuffled $r=0.341$; $p<1/200$). $T_i$
and $r_{500}$ are therefore not noise-dominated: they carry cluster-specific
information correlated with the CC-derived $H_\mathrm{CC}(z)$ values
\emph{through this specific quadrupole formula}. This disfavors explanation
(b) and tilts toward (a) --- the quadrupole functional form itself, not
proxy noise, drives the underperformance. This does not validate the fitted
$\beta_d,\beta_q$ coefficients (still phenomenological, \S\ref{ssec:beta}),
nor does it explain \emph{why} this particular combination of $k_A,
r_{500}, D$ underperforms mass-only scaling for real clusters --- both
remain open.
\reconstruction

\subsection{H$_\mathrm{FLRW}$ Baseline Provenance}
\label{ssec:hflrw}

To assess whether the tabulated $H_\mathrm{FLRW}(z)$ values correspond to
standard $\Lambda$CDM, we fit the 11 Table~A1 redshifts with three candidates:
(i)~Planck 2018 flat $\Lambda$CDM ($H_0 = 67.4$, $\Omega_m = 0.315$);
(ii)~a power law $H(z) = A(1+z)^p$; and (iii)~a best-fit CC-data flat $\Lambda$CDM.
The Planck model yields MAE $= 120$~km\,s$^{-1}$\,Mpc$^{-1}$ --- a failure by
any standard.  The power-law gives $A = 54.07$, $p = 0.884$,
MAE $= 5.04$~km\,s$^{-1}$\,Mpc$^{-1}$ (best single-model fit).
The CC-data $\Lambda$CDM fit converges to $H_0 = 68.78$, $\Omega_m = 0.316$,
$\chi^2/\mathrm{dof} = 0.93$.
\reconstruction

These results establish that TJB's $H_\mathrm{FLRW}$ baseline is not
Planck 2018 $\Lambda$CDM.  All MULTING vs.\ $\Lambda$CDM comparisons
in this paper use the power-law proxy as the H$_\mathrm{FLRW}$ reference
pending author confirmation of the intended convention.

\textbf{Effect on $\varepsilon(z)$ claims (EXP-G).}
To quantify the impact of the baseline choice on the MULTING expansion claims,
we compute $\varepsilon_\mathrm{Planck}(z) \equiv (H_\mathrm{MULT}/H_\mathrm{Planck})^2 - 1$
using the Planck 2018 flat $\Lambda$CDM baseline ($H_0 = 67.4$, $\Omega_m = 0.315$,
arXiv:1807.06209 Table~2).
The $H_\mathrm{FLRW}$ values fall progressively below Planck: $-1.9\%$ at $z=0.06$,
$-10.6\%$ at $z=0.40$, $-15.0\%$ at $z=0.65$, $-20.7\%$ at $z=1.00$,
and $-64.1\%$ at $z=8.50$.
As a consequence, $\varepsilon_\mathrm{Planck}(z)$ changes sign near $z \approx 0.33$:
\begin{itemize}
  \item $z < 0.33$: $\varepsilon_\mathrm{Planck} > 0$ --- MULTING predicts \emph{more}
        expansion than Planck $\Lambda$CDM.  This low-$z$ excess is
        baseline-independent: it survives rebasing from the power-law proxy
        to true Planck $\Lambda$CDM (e.g.\ $\varepsilon_\mathrm{Planck}(0) =
        \varepsilon_\mathrm{FLRW}(0) = +0.113$ at $z=0$, where both baselines
        share the same $H_0$).
  \item $z > 0.33$: $\varepsilon_\mathrm{Planck} < 0$ --- MULTING predicts
        \emph{less} expansion than Planck $\Lambda$CDM.
        At $z = 0.65$: $\varepsilon_\mathrm{FLRW} = +0.213$ vs
        $\varepsilon_\mathrm{Planck} = -0.124$ (sign reversal of $0.337$).
        The positive $\varepsilon$ in Table~A1 at $z \ge 0.40$ reflects
        the below-Planck $H_\mathrm{FLRW}$ baseline, not an independent signal.
        Claims based on $\varepsilon(z)$ for $z > 0.4$ should be interpreted
        relative to the stated baseline.
\end{itemize}
\reconstruction

\subsection{Physical proxy for the dipole component}
\label{ssec:fsig8}

The $\Lambda$CDM model prediction for the growth rate observable,
$f\sigma_8(z) = \Omega_m(z)^{0.55}\,\sigma_8\,D(z)$
(Carroll--Press--Turner growth factor), is non-monotone: it rises from
$f\sigma_8 \approx 0.449$ at $z=0.06$ to a peak of $\approx 0.601$ at $z \approx 0.65$,
then falls to $\approx 0.259$ at $z=5$.
This bell-shaped curve correlates with $\varepsilon(z)$ at Pearson
$r = 0.851$ across $z \le 5$ ($n=10$ Table~A1 points).
\tableAsubject

\emph{Subject of this correlation (added 2026-08-04).}
The $\varepsilon(z)$ values entering this and every following number in
this subsection are derived from Table~A1. Finding~B0 (2026-08-03,
\texttt{experiments/20260803-bridge/FINDING\_table\_a1\_provenance.md})
established from the author's own caption that each Table~A1 column,
$H_\mathrm{MULT}$ included, was requested from an online AI service
rather than computed from the force law, and that the material available
to us does not contain a deterministic, published, reproducible operator
$F \to H$ that produced it.
The shape correlation reported here is therefore a statement about that
AI response, not about MULTING. It is not thereby wrong --- the
arithmetic stands --- but its subject is not the theory, and no number in
this subsection may be quoted as a MULTING prediction in either
direction. A bridge test of the force law must target the
cosmic-chronometer compilation instead (\S\ref{sec:hz}).
The correlation captures a genuine \emph{shape match}: both are bell-peaked at
$z \sim 0.4$--$0.65$ ($r = 0.903$ on the rising half $z \le 0.65$;
$r = 0.958$ on the falling half $z \ge 0.65$).
Statistical significance: $p = 0.0018$ (two-tailed $t$-test, $df=8$);
95\% CI $= [0.48,\,0.96]$ (Fisher $z$-transform).
The correlation is robust across cosmological parameters:
$r \in [0.82, 0.87]$ on a $4\times4$ grid of
$\Omega_m \in [0.28, 0.32]$ and $\sigma_8 \in [0.77, 0.83]$.
\reconstruction

When tested against $n=11$ observed RSD $f\sigma_8$ measurements
(Beutler+2012, Blake+2012 WiggleZ, Samushia+2012 SDSS-LRG,
Howlett+2015 BOSS, de la Torre+2013 VIPERS, Okumura+2016 FastSound;
$z \in [0.07, 1.40]$; compiled from Kazantzidis \&\ Perivolaropoulos 2018,
arXiv:1803.01368, gold dataset), the correlation drops to $r = -0.084$
(EXP-D, [VERIFIED-REAL]), not significant ($p = 0.81$, $n=11$).
At the same $z$ values, the $\Lambda$CDM model gives $r = 0.778$ with $\varepsilon(z)$,
confirming that the gap is not a failure of the $\varepsilon$ shape but a
consequence of the $\sigma_8$ tension: observed $f\sigma_8$ values are
nearly flat ($\approx 0.38$--$0.51$), whereas the $\Lambda$CDM model
predicts a bell-peaked curve. \reconstruction
This absence of correlation with observed data may itself be informative:
both the MULTING $\varepsilon(z)$ deviation and the $\sigma_8$ tension
are largest near $z \sim 0.4$--$0.65$, suggesting a possible common
physical origin (open question, beyond BETA-1 HOLD).

This is consistent with MULTING's dipole term $F_d \propto k_A/r_P^2$
tracking the linear growth amplitude, pending $\beta_d$ derivation
(§\ref{ssec:beta}, open question Q2).

\textbf{Observational anchor caveat.}
The model $f\sigma_8(z)$ used in the correlation is a $\Lambda$CDM theoretical
prediction; observed RSD $f\sigma_8$ exist only up to $z \approx 2.33$
(DESI Lyman-$\alpha$, 2025, not yet validated for cosmological inference)
and $z \le 1.40$ for reliable galaxy-survey RSD.
All $f\sigma_8$ values at $z > 2$ in this analysis are model predictions,
not independent measurements.
\reconstruction

The full 11-point correlation yields $r = 0.775$, with the $z = 8.50$
point ($\varepsilon = 0.101 > \varepsilon_{z=5} = 0.048$, ratio $2.1\times$)
constituting an upward secondary bump.
This secondary feature at high redshift is consistent with the JWST-observed
excess of massive galaxies at $z \gtrsim 8$~\cite{tkachev2023}, which has been
attributed to a secondary enhancement in the dark-matter perturbation power
spectrum; the present framework offers an alternative geometric explanation via
the quadrupole term $F_q$.

\textbf{T-IDM (two-component structure, by physics-motivated exhaustion).}
Among 8 physics-motivated single-component families tested (virial $k_A(z)$,
merger-rate, log-Gaussian, Lognormal, $D^n H^m$ power products, and
merger-epoch $r_P$), no formula achieves Pearson $r > 0.85$ on all 11
$\varepsilon(z)$ points; the empirical ceiling is $r = 0.820$ for the
unconstrained $D^{5.5} H^{2.75}$ power product.
As an adversarial check, a 4-parameter Pad\'{e} rational function
$(a + bx)/(1 + cx + dx^2)$ with $x = 1+z$ achieves $r = 0.927$
on all 11 points~[\textit{OUR\_RECONSTRUCTION, verified numerically}].
However, the Pad\'{e} form carries no physical interpretation for the
secondary bump: its rational structure arises from fitting
degrees of freedom ($n/p = 2.75$), not from a physical mechanism.
This confirms that the 11-point $\varepsilon(z)$ curve requires a
non-monotone rational shape — consistent with MULTING's two-term
decomposition $F_\mathrm{oP} = F_m - F_d + F_q$, which provides a
physical identification of the two structural components (null-result
series NR-001 through NR-008, see supplementary).
\reconstruction

% ──────────────────────────────────────────────
\section{IDM Analysis}
\label{sec:idm}
% ──────────────────────────────────────────────

\subsection{Integer Isomer Count: $\chi^2$ Test}
\label{ssec:chi2}

The IDM prediction $\omega_\mathrm{DM} = N \omega_b$ (equal baryon asymmetry per
isomer) requires, from Planck 2018~\cite{planck2018}: $N_\mathrm{opt} =
\omega_\mathrm{DM}/\omega_b = 0.12011/0.022383 = 5.366$.

The nearest integer $N=5$ deviates by $5.79\sigma$. We test the point
hypothesis $H_0\!:\,\omega_\mathrm{DM} = N\omega_b$ directly via the exact
score-type statistic (variance evaluated at $N=5$, the null value, rather
than at the observed ratio): the linear combination
$T \equiv \omega_\mathrm{DM} - N\omega_b$ is exactly Gaussian under $H_0$
(no approximation is needed, unlike a delta-method expansion of the ratio
itself), with
\begin{equation}
\sigma_T = \sqrt{\sigma_{\omega_\mathrm{DM}}^2 + N^2\sigma_{\omega_b}^2}
= \sqrt{0.0012^2 + 5^2\times 0.00015^2} = 0.001415,
\end{equation}
yielding $T/\sigma_T = 5.79\sigma$ ($\chi^2 = 33.5$,
assuming equal baryon asymmetry per isomer, $f \equiv 1$;
see caveat below). A Wald-type delta-method estimate on the ratio
$N_\mathrm{opt} = \omega_\mathrm{DM}/\omega_b = 5.366$ itself (variance
evaluated at the \emph{observed} ratio) gives a close but distinct
$5.67\sigma$; we adopt the score-type statistic above as the point-null
test, since it requires no linearization and its variance is evaluated
under $H_0$, as a hypothesis test for a specific integer requires.
\reconstruction

As an independent cross-check, the DESI DR1+CMB combination~\cite{desi2024vi}
gives $\omega_b = 0.02248 \pm 0.00013$ and $\Omega_m h^2 = 0.1418 \pm 0.0028$,
yielding $N_\mathrm{opt}^\mathrm{DESI} = 5.31 \pm 0.13$,
consistent with the Planck value at $0.46\sigma$ (note: $\omega_\mathrm{cdm}$ is
a derived quantity from the DESI+CMB joint fit; DESI BAO alone does not constrain
$\omega_b$ independently~\cite{desi2024vi}).
\reconstruction

\textbf{Assumption caveat ($f \equiv 1$).}
The $5.79\sigma$ exclusion of integer $N=5$ is conditional on all dark isomers
sharing the same baryon asymmetry per comoving volume as ordinary matter
($f \equiv \eta_\mathrm{dark}/\eta_\mathrm{SM} = 1$).
If the dark-sector baryon asymmetry differs by an $O(1)$ factor $f \neq 1$,
then $\omega_\mathrm{DM} = N\,f\,\omega_b$ and any integer $N$ maps to
a physically admissible $f$: for $N=6$, one obtains $f = 0.894$, a mild
$\sim\!11\%$ deficit requiring no exotic physics.
The analysis therefore constrains $N_\mathrm{opt}$ as a ratio assuming $f=1$;
distinguishing integer $N$ values requires independent constraints on
$\eta_\mathrm{dark}$ (open question, beyond current scope).
\reconstruction

\textbf{Mass-spectrum resolution (EXP-N).}
A second resolution arises if the five IDM isomers carry unequal masses.
Under equal baryon number per isomer ($f \equiv 1$), the IDM prediction
becomes $\omega_\mathrm{DM}/\omega_b = \sum_i (m_i/m_p)$.  The Planck 2018
central value $N_\mathrm{opt} = 5.369$ is reproduced exactly by five isomers
with common mass $\bar{m} = 1.074\,m_p$ --- a $7.4\%$ enhancement over the
proton mass, or by any five-isomer spectrum whose masses sum to
$5.369\,m_p$~\cite{rosa2022}.  Such a mass splitting is within the
$O(1)$ range expected for a dark QCD analogue; it does not require exotic
physics.  Distinguishing this scenario from $f \neq 1$ requires
independent measurement of the dark-sector baryon asymmetry.
\reconstruction

\subsection{Hot Dark Matter Problem}
\label{ssec:hdm}

In a mirror-sector IDM, dark neutrinos acquire masses comparable to their
SM counterparts ($m_{\nu,\mathrm{dark}} \sim m_\nu$) and decouple
relativistically.  Their free-streaming length $\lambda_\mathrm{fs}$
erases density fluctuations below cluster scales, inconsistent with
the observed matter power spectrum and Lyman-$\alpha$ forest
($\lambda_\mathrm{fs} \gtrsim 10$~Mpc for $m_\nu \lesssim 0.1$~eV).
This makes mirror-sector IDM a hot dark matter (HDM) candidate at the
neutrino mass scale, ruled out by structure formation.
\reconstruction

A gravitational-only coupling (dark sector thermally isolated from the SM
bath) suppresses this problem by allowing $T_\mathrm{dark} \ll T_\mathrm{SM}$,
but eliminates the dark baryogenesis mechanism that generates the dark baryon
asymmetry required for IDM.  This tension (E4) is internal to the framework
and does not have a resolution within the preprint.

\subsection{Self-Interaction Cross Section}
\label{ssec:sigma}

The Bullet Cluster merger constrains the dark matter self-interaction
cross section to $\sigma/m < 1.25$~cm$^2$/g at 68\% confidence~\cite{markevitch2006}.

A geometric estimate for the dark-nucleon self-interaction
gives $\sigma_\mathrm{nuc}/m \approx \pi r_n^2/m_p \approx 0.027$~cm$^2$/g
(using $r_n = 0.87$~fm, $m_p = 938.3$~MeV/$c^2$), which is within the
Bullet Cluster limit by a factor of $\approx 46$.
However, if dark isomers contain dark atoms (dark proton $+$ dark electron),
the atomic cross section at cluster velocities
($v_\mathrm{rel} \sim 10^3$~km$/$s, Bohr radius $a_0 = 0.529$~\AA)
gives $\sigma_\mathrm{atom}/m \approx \pi a_0^2 / m_p \approx 5 \times
10^7$~cm$^2$/g, exceeding the limit by seven orders of magnitude.
For bound dark nuclei with $A \ge 1$, the cross section scales as
$\sigma \propto A^{2/3}$ while $m \propto A$, giving
$\sigma/m \propto A^{-1/3}$, which falls below the 1.25 limit for $A \ge 1$.
\reconstruction

This analysis implies that IDM dark matter must consist of dark nuclei
(not free dark protons or dark atoms) to satisfy the Bullet Cluster constraint.
A mirror-sector IDM with dark electromagnetism can form dark atoms (violating
the limit); a gravitational-only coupling eliminates dark electromagnetism
and renders the nucleon-level cross section consistent, but again requires
abandoning the dark baryogenesis mechanism (see also §\ref{ssec:neff}).
\reconstruction

\subsection{Effective Relativistic Species ($N_\mathrm{eff}$)}
\label{ssec:neff}

Three regimes must be distinguished; they differ by tens of orders of magnitude
in $\Delta N_\mathrm{eff}$ (EXP-O \reconstruction).

\textbf{(i) Thermalized mirror sector.}
If the dark sector was in full thermal equilibrium with the SM bath
($T_\mathrm{dark} = T_\mathrm{SM}$ at BBN), each mirror-sector isomer
contributes relativistic degrees of freedom.  For a minimal mirror sector
(dark photons + dark neutrinos) $\Delta N_\mathrm{eff} \approx 22$;
including dark quarks and gluons yields $\Delta N_\mathrm{eff} \approx 37$--81;
for the full mirror SM ($g_\mathrm{eff} = 106.75$), $\Delta N_\mathrm{eff} \approx 61$
(converting $g_*$ into neutrino-equivalent units, $\Delta N_\mathrm{eff} = g_*/[({7}/{8})\cdot 2]$
at $T_\mathrm{dark}=T_\nu$; quoting $g_*$ itself as $\Delta N_\mathrm{eff}$ conflates the two).
The Planck 2018 bound $\Delta N_\mathrm{eff} < 0.284$ (95\% CL,
$N_\mathrm{eff} = 2.99 \pm 0.17$) is exceeded by a factor of $\sim$80--215
in all three sub-cases --- qualitatively decisive; we quote the factor rather
than a Gaussian-tail $\sigma$-equivalent, which is not meaningful this far
outside the likelihood's validated range.  This is the regime implicitly
assumed by the earlier ``$>100\sigma$'' characterization.

\textbf{(ii) Early-decoupling hidden sector.}
If the dark sector decoupled from the SM at temperature $T_\mathrm{dec}$,
subsequent SM entropy production dilutes the dark temperature.  The ratio today
satisfies $\xi \equiv T_\mathrm{dark}/T_{\nu} = (10.75/g_{*S}(T_\mathrm{dec}))^{1/3}$.
For $T_\mathrm{dec} \gtrsim 100$~GeV, $g_{*S} = 106.75$, $\xi = 0.465$, and
$\Delta N_\mathrm{eff} = g_\mathrm{eff,dark}\,\xi^4 \approx 0.047\,g_\mathrm{eff,dark}$:
Planck 2018 allows $g_\mathrm{dark} \lesssim 6$ at 95\% CL~\cite{chacko2015}.

\textbf{(iii) Gravitational-only production (IDM as stated).}
For a sector that \emph{never} thermalizes with the SM, production is via
gravitational particle production during reheating.  The effective temperature
scales as $T_\mathrm{dark}/T_\mathrm{SM} \sim (T_\mathrm{reh}/M_\mathrm{Pl})^{4/3}$
\cite{clery2021}, giving $T_\mathrm{dark}/T_\mathrm{SM} \lesssim 10^{-11}$ for
$T_\mathrm{reh} \lesssim 10^{10}$~GeV.  The resulting
$\Delta N_\mathrm{eff} \sim g_\mathrm{dark}(T_\mathrm{dark}/T_\nu)^4 \lesssim 10^{-40}$
is entirely negligible~\cite{clery2021,farina2018}.
The $>100\sigma$ exclusion therefore applies exclusively to the thermalized
mirror-sector interpretation (\textbf{i}), not to gravitational-only IDM.
\reconstruction

The binding tension for IDM is thus \emph{not} $\Delta N_\mathrm{eff}$
(trivially satisfied in regime~iii), but the internal conflict that
grav-only coupling removes the dark baryogenesis mechanism IDM requires
to predict its dark abundance (E4 tension, §\ref{ssec:hdm}).
Quantitatively: matching $\omega_\mathrm{DM}/\omega_b = 5.369$ via five isomer
sectors requires per-sector number densities equal to the baryonic one to
$\sim$1.2\% ($n_i/n_b = 1.000 \pm 0.012$). A full thermal equilibrium
dense enough to fix $n_i \approx n_b$ over-produces $\Delta N_\mathrm{eff}$ by
the factors above, and the non-thermal (grav-only) history that satisfies
$\Delta N_\mathrm{eff}$ leaves the 5:1 ratio unexplained --- but these are not
the only two options: published cogenesis mechanisms (e.g.\ a common
``reheaton'' decaying into several sectors, transferring an asymmetry while
reheating them only to $T_\mathrm{dark}/T_\mathrm{SM}\lesssim1$, never fully
thermalizing with the SM) can in principle satisfy $\Delta N_\mathrm{eff}$
\emph{and} set definite per-sector densities in the same event
\cite{easa2022}. What no surveyed mechanism --- thermal, non-thermal, or
transfer-based --- produces is the specific structure IDM's isomer count
needs: \emph{five identical} sectors each independently landing at
$n_i/n_b\approx1$; known multi-sector constructions reach a similar total
ratio through internal structure other than five equal densities (e.g.\
differential decoupling across many \emph{unequal} sectors, or a few
components equalized by conversions). The 5:1 match therefore remains a
posited coincidence rather than a derived consequence of any known cogenesis
history, independent of the specific $\Delta N_\mathrm{eff}$ escape invoked.
\reconstruction

% ──────────────────────────────────────────────
\section{Discussion}
\label{sec:discuss}
% ──────────────────────────────────────────────

\subsection{What Survives: H4 Empirical Relations}

Two results from this assessment are independent of the MULTING/IDM
cosmological framework and survive even if the framework itself is not
confirmed.

First, the fermion-gravity link (Eq.~32) holds at $1.0\sigma$ using PDG 2024
constants.  The relation $(4/3)(m_\tau/m_e)^{12} = \alpha_\mathrm{EM}/\alpha_G$
connects the strong gravitational coupling to the electromagnetic fine-structure
constant through lepton masses.  Its mechanism is unknown, but its numerical
precision is not easily dismissed as coincidence.

Second, the boson mass ratio $m_W^2:m_Z^2:m_H^2 \approx 7:9:17$ holds to
$<0.05\%$ in mass units across all three bosons in a least-squares fit, but
its error-weighted significance is anchor-dependent and dominated by the
$W$-boson residual ($\approx 3.8\sigma$ when $Z$-anchored, vs $1.14\sigma$ for
the Higgs); we therefore present it as a precise but not yet error-consistent
regularity. These ratios may encode a symmetry group that has yet to be identified.

Both results are independently testable as new PDG data accumulates (the
$m_\tau$ uncertainty currently dominates Eq.~32) and constitute H4 of the
original assessment framework.

\textbf{Uniqueness of $n=12$ in Eq.~(32).}
A power scan over integer exponents $n$ confirms that $n=12$ is the
\emph{unique} integer for which the relation
$(4/3)(m_\tau/m_e)^n = \alpha_\mathrm{EM}/\alpha_G$ holds within any
reasonable tolerance: adjacent integers give deviations of $\gtrsim 10^5\%$
($n=11$: $-99.97\%$; $n=13$: $+3.5\times10^5\%$). The best-fit
continuous exponent is $n_\mathrm{best} = 12.000$, confirming that $n=12$
is not an arbitrary choice but the unique integer where the ratio
$(m_\tau/m_e)^n$ crosses the target $\alpha_\mathrm{EM}/\alpha_G$.
A comprehensive look-elsewhere scan over $5{,}640$ combinations
(all reduced rational prefactors $p/q$ with $p,q\le 12$, integer exponents
$1\le n\le 20$, and the three lepton mass pairs $\tau/e$, $\mu/e$, $\tau/\mu$,
all targeting the same RHS $\alpha_\mathrm{EM}/\alpha_G$)
yields a \emph{single unique reduced-form match} within $1\%$, namely
$(4/3,\,n{=}12,\,m_\tau/m_e)$; the $\mu/e$ and $\tau/\mu$ pairs have zero
matches at any threshold.
The raw count of three hits ($4/3$, $8/6$, $12/9$) all reduce to the same
unique fraction, so the uniqueness is of the prefactor and mass-pair
combination, not merely of the exponent.
Note that for any fixed mass pair and target, the large ratio
$\alpha_\mathrm{EM}/\alpha_G \approx 4{\times}10^{42}$ guarantees exactly
one integer exponent will be nearest the target (adjacent exponents deviate
by $\gtrsim 10^5\%$), so the non-trivial claim is the uniqueness of the
\emph{prefactor} $4/3$ and \emph{mass pair} $\tau/e$, both of which
survive the comprehensive scan.
An extended scan over $83{,}160$ combinations --- adding all quark masses,
the proton and neutron (PDG 2024 values), prefactors $p/q$ with $p,q\le 10$,
and exponents $1\le n\le 24$ --- confirms that Eq.~(\ref{eq:c9}) remains
rank~$\#1$ with a $0.014\%$ deviation; the second-best formula
$(7/5)(m_t/m_s)^{13}$ deviates by $0.30\%$, a factor of~22 larger.
Within the $0.1\%$ threshold, Eq.~(\ref{eq:c9}) is the sole match
(empirical $p < 2\times10^{-5}$).
\reconstruction
\reconstruction

\textbf{Algebraic origins of the coefficients.}
Independently of any physical model, the three integers appearing in
Eq.~(\ref{eq:c9}) --- the prefactor $4/3 = 36/27$, the exponent $12$, and
the implicit product $36$ --- correspond to dimensions of specific algebraic
objects in the exceptional-Lie-algebra chain $\mathrm{G}_2 \subset \mathrm{F}_4$:
(i) $12 = |\text{roots}(\mathrm{G}_2)|$, the total number of non-zero roots of
the Lie algebra $\mathrm{G}_2 = \mathrm{Aut}(\mathbb{O})$, where $\mathbb{O}$
denotes the octonions \cite{slansky1981} --- for a rank-2 root system this
equals the order of the Weyl group $|W(\mathrm{G}_2)|$ by construction, so the
two are not independent observations of the same integer, only one; and
independently, $12$ is the highest degree among the four fundamental Casimir
invariants of $\mathrm{F}_4$ (degrees $2, 6, 8, 12$) \cite{bincer1993f4casimir}
--- two independent algebraic routes to the same integer, not three. (The
Casimir-degree property is also not unique to $\mathrm{F}_4$: $\mathrm{E}_6$,
$\mathrm{E}_7$, and $\mathrm{E}_8$ each likewise include $12$ among their own
Casimir/exponent degrees, so this route does not by itself single out
$\mathrm{F}_4$.)
(ii) $36 = \dim(\mathrm{SO}(9))$, and $27 = \dim(J_3(\mathbb{O}))$, the dimension
of the Albert algebra (exceptional Jordan algebra), giving $4/3 = 36/27$ ---
one of several comparable-weight ratios available within the same $\mathrm{F}_4$
decomposition (e.g. $52/36$, $52/27$, $16/9$), selected here because it matches
the target value.
A null test over 2256 pairs from a pool of 48 algebraic integers (squares and
integer roots of dimensions of classical and exceptional Lie algebras)
confirms that finding both the ratio $4/3$ and the subset $\{36,27,12\}$ by
chance has probability $p\approx 1/37000$.
These connections are empirical observations, not derivations: they identify
the algebraic objects that \emph{count} the coefficients but do not explain
why those objects appear in the ratio $\alpha_\mathrm{EM}/\alpha_G$.
The $\mathrm{G}_2$ route is natural because $\mathrm{G}_2 = \mathrm{Aut}(\mathbb{O})$
and octonionic geometry underlies several approaches to Standard Model
structure~\cite{baezhuerta2010, todorov2018}; the specific connection
to $\alpha_G$ remains an open question.
A quark analogue --- replacing $m_\tau/m_e$ with quark mass ratios and the
exponent 12 with $\mathrm{E}_8$-motivated integers --- was tested and not found:
no quark pair satisfies an analogue of Eq.~(\ref{eq:c9}) within $5\%$ for
any exponent in the relevant range, confirming that the relation is
\emph{lepton-specific}.
\reconstruction

\textbf{$m_\tau$ value for an exact relation.}
If Eq.~(32) is taken as exact, it predicts
$m_\tau^\mathrm{exact} = 1776.840$~MeV,
which differs from the PDG 2024 value ($1776.93 \pm 0.09$~MeV) by
$-0.090$~MeV $= 1.00\sigma$.
Belle~II has already measured $m_\tau$ once, in 2023~\cite{belleii2023taumass}:
$1777.09 \pm 0.08\,\mathrm{(stat)} \pm 0.11\,\mathrm{(syst)}$~MeV, a
$1.84\sigma$ tension with the exact-relation value above (added 2026-09-03;
see \S\ref{ssec:eq32} for the full computation) --- ambiguous, not
decisive. Future Belle~II Run~2 measurements targeting
$\sigma_{m_\tau} < 0.01$~MeV will sharpen this toward a definitive test of
whether Eq.~(32) is exact or approximate.
\reconstruction

\textbf{Joint Koide--Eq.~(32) constraint.}
The Koide formula
$(m_e + m_\mu + m_\tau)/(\sqrt{m_e}+\sqrt{m_\mu}+\sqrt{m_\tau})^2 = 2/3$
holds at $0.00033\%$ ($0.43\sigma$, PDG 2024) and is symmetric under all six
permutations of $\{m_e, m_\mu, m_\tau\}$ (i.e., $S_3$-invariant), involving
all three lepton generations.
Eq.~(\ref{eq:c9}) is distinct: it connects only the first and third
generations ($m_\tau/m_e$) to the gravitational coupling.
Inverting each relation for $m_\tau$ (the Koide prediction depends only on
$m_e,m_\mu$ and is unaffected by the PDG-2024 tau-mass correction above):
\begin{align}
  m_\tau^\mathrm{Koide} &= 1776.969~\mathrm{MeV} \quad (+0.43\sigma\text{ from PDG}) ,\\
  m_\tau^\mathrm{Eq.32} &= 1776.840~\mathrm{MeV} \quad (-1.00\sigma\text{ from PDG}) .
\end{align}
The two ``exact'' predictions differ by $+0.129$~MeV ($+1.43\sigma$) ---
marginally compatible.  If both relations are simultaneously exact,
$m_\tau$ must lie within the window $[1776.840, 1776.969]$~MeV, which is
narrower than the current PDG uncertainty ($\pm 0.09$~MeV, full $1\sigma$
band $0.18$~MeV) and would be resolved by the Belle~II target precision.
Under this joint constraint, Newton's constant $G$ is determined from
$\alpha_\mathrm{EM}$, $m_e$, and $m_\mu$ alone, without $m_\tau$ as a free
parameter --- a non-trivial cross-constraint whose mechanism is unknown.
\reconstruction

\subsection{What Requires New Physics}

Three open questions require new physics or author disclosure to resolve.

\textit{MULTING Lagrangian (G3).}
No Lagrangian or action principle is provided for the multipole force
$F_\mathrm{oP}$.  Without a Lagrangian, the theory cannot be connected to
standard field-theoretic tools (renormalization, symmetry analysis, gravitational
waves) and Lorentz covariance cannot be verified.

\textit{IDM symmetry group (G4).}
The six IDM isomers are postulated by analogy with the SM, but the symmetry
group that distinguishes exactly six from other values of $N$ is not derived.
The $5.79\sigma$ tension with Planck 2018 ($N_\mathrm{opt} = 5.366$) may be
resolvable if the symmetry group predicts non-equal isomer masses or
baryon asymmetries.

\textbf{$S_3$ as a candidate group (hypothesis).}
The symmetric group $S_3$ (permutations of 3 objects, order 6) arises
independently in two contexts that converge on IDM's isomer count.
First, the Koide formula
$(m_e+m_\mu+m_\tau)/(\sqrt{m_e}+\sqrt{m_\mu}+\sqrt{m_\tau})^2 = 2/3$
is invariant under all six permutations of $\{m_e, m_\mu, m_\tau\}$
(i.e., $S_3$-symmetric by construction, verified independently~\cite{pdg2024}).
Second, under the regular (Cayley) representation of $S_3$ on itself,
the elements split canonically into $\{e\} \cup (S_3 \setminus \{e\})$:
one identity element (SESI, ordinary matter) and five non-identity elements
(MESI, dark matter), giving DM:OM $= 5{:}1$ as a group-theoretic consequence
rather than a counting postulate.
Additionally, the standard model's non-commutative geometry (Connes--Lott)
assigns $|S_3| = 6$ summands to the doubled algebra $\mathcal{A}_F \otimes
\mathcal{A}_F^\mathrm{op}$ and KO-dimension $= 6$, offering a third
independent path to the same integer.
This observation does not constitute a derivation; it identifies $S_3$ as a
candidate symmetry whose structure forbids $N=4$ or $N=7$ and predicts
the specific $1+5$ split without adjustable parameters.
Confirming this identification requires showing that IDM isomer quantum numbers
transform under the defining representation of $S_3$, which we record as an
open hypothesis.
\reconstruction

\textbf{$\Delta N_\mathrm{eff}$ constraint on IDM dark neutrinos (hypothesis).}
If each MESI isomer mirrors the SM, it contains three generations of dark
neutrinos, giving $N_{\nu,\mathrm{dark}} = 5 \times 3 = 15$ light fermionic
species.  When the dark sector is in thermal equilibrium with the SM at
temperature $T_\mathrm{dec}$ and subsequently decouples, SM entropy injections
(e.g., $e^+e^-$ annihilation, QCD transition) heat the photon bath relative
to the dark sector.  The residual dark-sector temperature ratio is
$\xi \equiv T_\mathrm{dark}/T_\gamma = (g_{*S}^\mathrm{today}/g_{*S}(T_\mathrm{dec}))^{1/3}$,
and the contribution to the effective number of relativistic species is
\begin{equation}
  \Delta N_\mathrm{eff} = N_{\nu,\mathrm{dark}}\,\xi^4\,\left(\tfrac{11}{4}\right)^{4/3}.
  \label{eq:dneff}
\end{equation}
At the SM maximum $g_{*S}(T_\mathrm{dec}) = 106.75$ (all SM particles in equilibrium
at $T > 200\,\mathrm{GeV}$), one finds $\xi = 0.332$ and
$\Delta N_\mathrm{eff} = 0.70$.
The Planck 2018 constraint $N_\mathrm{eff} = 2.99 \pm 0.17$ (68\% CL)
limits $\Delta N_\mathrm{eff} < 0.29$ (95\% CL), so 15 dark neutrinos are
\emph{ruled out} at the 2.5$\times$-excess level if the dark sector
thermalised within the SM.
The constraint requires either (A) dark-sector decoupling at $T_\mathrm{dec}$
where $g_{*S} > 210$, achievable in MSSM-like extensions with
$g_{*S}^\mathrm{MSSM} \approx 228$ at $T > 1\,\mathrm{TeV}$, or
(B) the absence of light dark neutrinos.
An alternative exit via massive dark neutrinos also fails: the minimum
mass for relativistic suppression implies $\Sigma m_{\nu,\mathrm{dark}} \gtrsim 1.3\,\mathrm{eV}$,
exceeding the Planck neutrino mass limit $\Sigma m_\nu < 0.12\,\mathrm{eV}$.
If the MSSM scenario (A) holds, the predicted $\Delta N_\mathrm{eff} \approx 0.26$
is above the Simons Observatory 2$\sigma$ detection threshold ($\sigma \approx 0.05$)
and will be tested by 2027.
\reconstruction

\textit{$\beta$ derivation from first principles (G2).}
The parameters $\beta_d$ and $\beta_q$ are fitted phenomenologically.
Their $5.8\times$ and $95\times$ spreads across AI services indicate
that the current procedure does not uniquely determine them.
A first-principles derivation linking $\beta$ to cluster microphysics
is the single most impactful open question for the MULTING H(z) test.

Note: Blanchet \& Le Tiec~\cite{blanchet2008, blanchet2009} demonstrated
that a dipolar dark-matter fluid is degenerate with $\Lambda$CDM at linear
perturbation order.  Our 5-fold cross-validation (§\ref{sec:hz}) confirms
this numerically: $\Delta\mathrm{AIC} = 0.74$, indistinguishable.
MULTING requires second-order predictions (non-geodesic cluster motions,
lensing anisotropy) to become observationally distinguishable.

\subsection{Forward-Looking Tests}
\label{sec:predictions}

The experiments reported here yield five falsifiable predictions, each
testable with facilities operating or under construction.
Table~\ref{tab:predictions} summarises the predictions, the required
observable, the instrument, and the approximate timeline.

\begin{table*}[tb]
\caption{Falsifiable predictions from this analysis (OUR\_RECONSTRUCTION).
All predictions assume currently favoured cosmological parameters;
updates in DESI DR2 or Planck PR4 may shift central values.}
\label{tab:predictions}
\begin{ruledtabular}
\begin{tabular}{p{5.5cm}p{3.5cm}p{2.5cm}l}
Prediction & Observable & Instrument & Timeline \\
\hline
$m_\tau^\mathrm{exact} = 1776.840$~MeV if Eq.~(32) exact ($-1.00\sigma$ from PDG; already $1.84\sigma$ from Belle~II's own 2023 measurement~\cite{belleii2023taumass}) &
$m_\tau$ to $\pm0.01$~MeV & Belle~II Run~2~\cite{belleii2019} & 2027 \\
$\Delta N_\mathrm{eff} \approx 0.26$ if IDM + MSSM decoupling ($g_{*S} \approx 228$) &
$\Delta N_\mathrm{eff}$ & Simons Obs. & 2026--2027 \\
$N_\mathrm{opt} = 5.31 \pm 0.05$--$0.09$ (3.3--5.9$\sigma$ from $N{=}5$) if $\sigma_{\omega_m}$ drops $\sqrt{2}$--$3\times$ from DR1 &
$\omega_m$, $\omega_b$ & DESI~DR2 & 2025--2026 \\
Non-geodesic cluster motions at $\gtrsim 2\sigma$ if $\beta_d \gtrsim 10^4$ &
$\sigma_v$, lensing anisotropy & Euclid~DR1 & 2026+ \\
$\varepsilon(z)$ secondary peak at $z \approx 8.5$ from $F_q$ if $\beta_q \gg \beta_d$ &
$f\sigma_8(z)$ at $z > 5$ & Future & 2030+ \\
\end{tabular}
\end{ruledtabular}
\end{table*}

The most time-critical test is the Simons Observatory $\Delta N_\mathrm{eff}$
measurement.
If IDM's mirror sector decoupled from the SM at $T_\mathrm{dec}$ where
$g_{*S}(T_\mathrm{dec}) \approx 228$ (MSSM regime), the predicted
$\Delta N_\mathrm{eff} \approx 0.26$ lies above the Simons Observatory
$2\sigma$ detection threshold ($\sigma_{\Delta N_\mathrm{eff}} \approx 0.05$)
and below the Planck 2018 $95\%$ CL limit ($0.29$).
A non-detection by Simons Observatory would require $g_{*S}(T_\mathrm{dec}) > 300$,
pushing IDM's thermal decoupling to scales above the MSSM.
\reconstruction

\subsection{Assessment Summary}

\begin{table*}[tb]
\caption{36-point assessment summary with candidate resolution paths
(OUR\_RECONSTRUCTION). Each gap is paired with a constructive next step;
paths marked ``author'' require the framework author's disclosure or
input and are not endorsed here as fixes.}
\label{tab:summary}
\begin{ruledtabular}
\begin{tabular}{p{2.7cm} p{4.2cm} p{2.1cm} p{4.9cm}}
\hline
Component & Status & Evidence & Path to resolution / next test \\
\hline
Eq.~31 (7:9:17 bosons) & $\triangle$ $m_W/m_H$ at $0.04\%$; $1.14\sigma$ (H), $3.8\sigma$ (W); trials-corrected $p = 0.0051 \pm 0.0002_\mathrm{MC}$ ($a,b,c \le 20$), but $\sigma_\mathrm{spec}$ spans $0.003$--$0.077$; isolation test $P=0.34$ (negative) & PDG 2024 $+$ 997-triple null & $W$-mass world average (LHC Run~3); $m_H$ precision also binds \\
Eq.~32 (fermion-gravity) & $\checkmark$ $1.0\sigma$ (PDG); $1.84\sigma$ (Belle~II 2023, independent, ambiguous); trials-corrected $p = 0.16$ & PDG 2024 $+$ 35\,672-expression null $+$ Belle~II 2023 direct check & Belle~II Run~2 $m_\tau$ (2027) sharpens exact vs.\ approximate \\
Eqs.~21--24 (fermion masses) & $\checkmark$ muon $0.47\%$, quarks $<0.31\%$ & PDG 2024 & confirmed asset \\
Inflaton $m_Z/3 = 30.4$~GeV & ? undetected & LEP/LHC null & precision/collider search \\
Cassini constraint & $\checkmark$ pass & $\varepsilon \sim 10^{-19}$ & satisfied \\
H(z) correlation & $\triangle$ $\Delta\mathrm{AIC}{=}{+}0.74$ & 27 Moresco CC & out-of-sample test needs the author's $F_{oP}\!\to\!H(z)$ bridge \\
$H_\mathrm{FLRW}$ provenance & $\triangle$ not Planck baseline & MAE $120$ vs $5$ & recompute on the Planck~2018 baseline; author to confirm intended convention \\
$\beta_d = 4.5$ observability & $\triangle$ negligible as tabulated & $k_A/M$ calc & first-principles $\beta$ (Lagrangian) or author disclosure; not recoverable by data calibration (non-identifiable) \\
$\beta$ provenance (BRAI) & $\triangle$ $R_B{=}15.9/24.1$ ($p{<}10^{-4}$) & 3 AI services & author's intended $\beta$ definition \\
IDM $N=5$ equal-mass & $\triangle$ $5.79\sigma$ if equal & Planck 2018 (EXP-N) & relax to unequal isomer masses (next row) \\
IDM $N=5$ unequal-mass & $\triangle$ $\bar{m}=1.074\,m_p$ & mass spectrum & sum matches $\omega_\mathrm{DM}/\omega_b$; adds a parameter, needs independent motivation \\
IDM cold-DM viability & $\triangle$ warm-DM free-streaming tension & thermal history & warm/mixed-DM or massive dark $\nu$; author to specify thermal history \\
$f\sigma_8$ dipole proxy & $\triangle$ $r{=}0.851$ ($n{=}10$, $z{\le}5$) --- \emph{subject is the Table~A1 AI response, not MULTING (B0)} & Table A1 $+$ $\Lambda$CDM & retest against cosmic chronometers, not Table A1 \\
Secondary $z{=}8.5$ bump & $\triangle$ structural --- \emph{same subject caveat (B0)} & Two-component $\varepsilon$ & quadrupole $\beta_q$ from author \\
MULTING Lagrangian & ? open & Future work & action-principle derivation (G3): the key open problem \\
6-isomer symmetry & ? open & Future work & $S_3$ representation test (order $6 = 1{+}5$) \\
\hline
\end{tabular}
\end{ruledtabular}
\end{table*}

% ──────────────────────────────────────────────
\section{Conclusions}
\label{sec:conclusions}
% ──────────────────────────────────────────────

We have conducted a systematic 36-point assessment of the MULTING/IDM
framework~\cite{buckholtz2026}, focusing on components that are independently
verifiable with public data and PDG constants.

Several results survive independent scrutiny without requiring the author's
undisclosed numerical procedure.
The fermion-gravity link (Eq.~32) holds at $1.0\sigma$ (PDG 2024),
and the boson mass ratio [7:9:17] holds at $<0.05\%$ in mass units
with the scale-free ratio $m_W/m_H = \sqrt{7/17}$ precise to $0.04\%$.
The fermion mass spectrum (Eqs.~21--24) reproduces the muon mass to $0.47\%$
and quark geometric means to $<0.31\%$.
These H4 results are non-trivial empirical regularities regardless of the
status of MULTING cosmology.

Their strength should not be overstated, and we quantify one of them rather
than assert it.  Eq.~32's $1.0\sigma$ measures experimental precision given the
formula; the search that produced the formula is not free.  Enumerating 170\,352
simple prefactor--power expressions, of which 35\,672 are statable to better than
$0.5\%$ after propagating their own mass uncertainties, and comparing against
4000 random targets of the same magnitude, gives
$p_\mathrm{look\text{-}elsewhere} = 0.16$ (\S\ref{ssec:eq32}).
Eq.~32 is twice better than the median such target and remains the sole statable
hit at its own precision, but its agreement is ordinary rather than
extraordinary.

The same treatment applied to [7:9:17] gives the opposite verdict, and the two
should not be quoted as though they carried equal weight.  That relation needs a
separate enumerator --- it fixes two scale-free ratios with a triple of small
integers rather than a prefactor times a power --- scored by the minimax residual
over a free scale, which is the paper's own figure of merit
(\S\ref{ssec:bosons}).  Among the 997 coprime triples $a<b<c\le 20$,
$(7,9,17)$ is the \emph{unique} best fit at $0.050\%$, against a median of
$0.64\%$ for random targets drawn within a factor 1.5 of the observed ratios:
\begin{equation}
  p_\mathrm{look\text{-}elsewhere}^{[7:9:17]} = 0.0051 \pm 0.0002
  \quad (a,b,c \le 20) ,
\end{equation}
from $K=233{,}726$ random targets.  The quoted uncertainty is binomial and is
not optional --- single runs of $\sim\!2\times10^{4}$ targets scatter over
$0.003$--$0.010$ --- but it is also \emph{not} the dominant uncertainty.  It
measures only $\sigma_\mathrm{MC}$; the specification uncertainty
$\sigma_\mathrm{spec}$, induced by the choice of grammar, spans
$0.003$--$0.077$ across the families scanned below, two orders of magnitude
larger.  No single number is therefore ``the'' significance of this relation,
and the figure above should never be quoted without its bound.
It remains the unique best triple out to $a,b,c \le 50$, and it wins in
$100\%$ of $3000$ pseudo-datasets drawn from the PDG uncertainties.  That second
figure must not be read as a look-elsewhere result: it asks whether measurement
noise of order $0.088\%$ can dislodge a triple already sitting $0.050\%$ away,
and the answer is almost tautologically no.  It establishes robustness of the
rank against experimental error and says nothing about how often a random target
would find so good a candidate.  Likewise, the grammars in which the rank
persists are largely nested --- integers $\le 20$ is contained in integers
$\le 50$ and in the rational families --- so persistence of the rank under
enlargement demonstrates that no competitor appears, which is weaker than
demonstrating that none was expected to.  Triples sharing a common divisor are
excluded as duplicates rather than competitors: $(14,18,34)$ states the same
proportion as $(7,9,17)$.  The ordering $a<b<c$ is forced by
$m_W<m_Z<m_H$ and carries no additional assignment freedom.  Working in mass
units rather than squared-mass units rescales every residual by exactly two and
therefore leaves the rank and $p$ unchanged.

The bound $N$ must be quoted with the number, and a grammar scan shows why.  We
repeated the identical measurement across twelve declared grammars --- integers
to several bounds, ratios of small integers, $7$- and $11$-smooth numbers,
squares, and arithmetic progressions --- canonicalising every candidate to its
unique coprime integer triple, since a scale-free proportion has one such form.
Among the seven grammars that can express $(7,9,17)$, all of which return the
same $0.050\%$ residual and are therefore scored against an identical threshold,
\begin{equation}
  p \;=\; (7.1\times10^{-6})\,|G|^{0.95},
  \qquad \text{median residual } 2\% ,
\end{equation}
with the integer and rational families \emph{interleaved} on that line rather
than occupying separate curves.  An exponent near unity is not a discovery: for
a nearest neighbour among $|G|$ candidates at roughly uniform density, the
rare-event probability of landing within a fixed radius grows as $|G|\,V_d(r)$,
so $p \propto |G|^{1}$ is precisely what a well-behaved enumerator should
return.  Its value here is diagnostic rather than physical --- it confirms that
$p$ is being set by candidate density and essentially nothing else.  Ratios of small integers are thus not a second
degree of freedom: canonicalised, they are a differently shaped subset of the
same triples and cost precisely what their count costs.  The trials factor is,
to within tens of percent, the size of the declared search and little else ---
$p$ rises to $0.077$ at $N=50$ for that reason.  It is by contrast stable
against changes in the null's sampling window.

Two further points follow, one of which weakens the result and one of which
sharpens it.  First, no principled cut on ``simple'' exists, so the bound is a
choice we make and the $p$ must always be quoted with it.  Second, ``simple'' is
not one axis: five of the twelve grammars cannot state the relation \emph{at
all}, among them $11$-smooth numbers below $200$, a space three times larger
than integers $\le 50$.  The triple is small by magnitude but not simple by
factorisation --- $17$ is prime --- nor by pattern, being neither an arithmetic
progression nor a ratio of squares.  A trials factor quoted without naming which
axis of simplicity it ranges over is incomplete.

A sharper test separates the two hypotheses that a bare $p$ conflates: that the
triple has some structure making it unusually close to the data ($H_1$), against
that it is simply the nearest neighbour in a sufficiently dense discrete space
($H_0$).  A genuinely special point should be not only \emph{close} but
\emph{isolated}, so we record for every target both the best residual $D_1$ and
the gap to the runner-up, $D_2-D_1$.  The result does not favour $H_1$.  For
$a,b,c\le 20$ the observed ratio is $D_2/D_1 = 14.7$, whereas random targets
that are equally close have a \emph{median} $D_2/D_1$ of $18.0$; the conditional
probability $P(D_2-D_1 \ge \text{observed} \mid D_1 \le \text{observed})$ is
$0.34$, and $0.19$ at $N=50$.  Conditional on being close, $(7,9,17)$ is no more
isolated than a typical random target --- at $N=20$, slightly less.  The
isolation dimension therefore contributes no evidence beyond closeness itself,
and closeness is exactly the quantity shown above to be set by the density of
the declared grammar.

\emph{These two $p$-values are not directly comparable as measures of evidential
strength.}  Eq.~32's grammar is one-dimensional and Eq.~31's is
two-dimensional, and the corrected tail probabilities were computed under
different pre-specified search spaces; harmonising the null spaces would be
required before the ratio of the two numbers meant anything.  The defensible
statement is the weaker one: under its own declared grammar each relation
carries a measured trials factor, and that factor is substantially smaller for
[7:9:17].  Both figures remain conditional on grammars chosen by us
\emph{after} the relations were known, which no amount of enumeration removes.

Finally, the $0.050\%$ residual sits below the $0.088\%$ that the current $m_H$
uncertainty can resolve.  This is a statement about what the data cannot do, not
about what they confirm: the measurements are \emph{compatible} with $7:9:17$,
but they lack the precision to distinguish it from a whole neighbourhood of
nearby proportions.  Compatibility, precision, and evidence for an exact
relation are three different claims, and only the first is currently supported.
Two consequences follow.  First, $m_H$ precision joins the $W$-mass combination
as a discriminant.  Second, the relation makes point predictions sharper than
present data, which future measurements will test out of sample even though the
present ones cannot: anchoring on $m_Z$ gives $m_H = 125.3254$~GeV
($+1.14\sigma$) and $m_W = 80.4199$~GeV ($+3.92\sigma$ against the world average
adopted here).  The $W$ prediction is thus already in tension, and no part of
the agreement reported above is out-of-sample --- all three masses entered the
selection of the triple.

A key methodological finding is that the H$_\mathrm{FLRW}$ baseline in
Table~A1 is inconsistent with standard Planck 2018 $\Lambda$CDM
(MAE $= 120$~km\,s$^{-1}$\,Mpc$^{-1}$); a power-law $H(z)=54.07(1+z)^{0.884}$
fits the same 11 points with MAE $= 5.04$~km\,s$^{-1}$\,Mpc$^{-1}$.
This means that comparisons of MULTING vs.\ ``standard cosmology'' must
specify which H$_\mathrm{FLRW}$ convention is intended.

The IDM cosmological predictions face two structural constraints, with a
third that depends on the coupling assumption.
First, $N=5$ \emph{equal-mass} isomers are excluded at $5.79\sigma$ by Planck
2018; however, five isomers with a common mass $\bar{m} = 1.074\,m_p$ exactly
reproduce the observed $\omega_\mathrm{DM}/\omega_b = 5.369$, converting the
exclusion into a mass-spectrum constraint (EXP-N).
Second, the $\Delta N_\mathrm{eff}$ bound depends critically on the dark-sector
coupling: a fully thermalized mirror sector gives $\Delta N_\mathrm{eff} \approx 61$,
exceeding the Planck bound by a factor of $\sim$215; a gravitational-only sector (IDM as stated) gives
$\Delta N_\mathrm{eff} \lesssim 10^{-40}$, entirely within Planck bounds (EXP-O).
The ``$>100\sigma$'' characterization applies only to the thermalized regime
and is a \emph{misclassification} when applied to grav-only IDM~\cite{clery2021}.
A massive-neutrino escape route for the thermalized case fails because it implies
$\Sigma m_{\nu,\mathrm{dark}} \geq 1.3$~eV $\gg$ the Planck mass limit
$0.12$~eV.
The minimal viable scenario requires BSM physics above the electroweak scale
($g_{*S}(T_\mathrm{dec}) > 210$); MSSM-like extensions predict
$\Delta N_\mathrm{eff} \approx 0.26$, within Planck bounds but detectable by
Simons Observatory by 2027 (Table~\ref{tab:predictions}).
Third, a gravitational-only coupling that would suppress $\Delta N_\mathrm{eff}$
eliminates the dark baryogenesis mechanism (E4 internal tension).

The MULTING H(z) test is currently limited by the $\beta$-rescaling gap:
the tabulated $\beta_d = 4.5$ gives a negligible dipole signal
($\varepsilon \sim 10^{-6}$), while the AI-service values
($\beta_d \sim 10^4$) are not derived from first principles.
The growth-rate observable $f\sigma_8(z)$ reproduces the shape of
$\varepsilon(z)$ at $r = 0.851$ for $z \le 5$ and identifies a secondary
structural feature at $z = 8.50$ consistent with the quadrupole term $F_q$
--- with the subject caveat of finding~B0: $\varepsilon(z)$ here is derived
from Table~A1, so both statements describe the AI service's response and
not MULTING.
\tableAsubject

The primary open questions are: a first-principles derivation of $\beta_d$
and $\beta_q$, a Lagrangian for MULTING, and the symmetry group for the
IDM isomer count.
A fourth, established after this draft's last substantive revision, is
prior to all of them: whether a MULTING calculation of $H(z)$ exists that
does not pass through an AI service.
If it does, that calculation --- not Table~A1 --- is the object to
reconstruct; if it does not, then building a bridge is new work rather
than reconstruction.
Pending these, no conclusions about the physical reality of the dipole
force can be drawn from Table~A1 in either direction.

\paragraph*{Code and data availability.}
All numerical checks in this paper are reproducible via
\texttt{scripts/verify\_all\_claims.py}, which runs in under five seconds
on a standard laptop and writes a machine-readable JSON report.
The harness and all experiment scripts are publicly available at\\
\url{https://github.com/sergeeey/buckholtz-idm-multing-study}\\
(branch \texttt{main}; commit hash recorded in the JSON report).
PDG 2024 constants are embedded in the script; no proprietary data are used.

\begin{acknowledgments}
S.B.\ thanks T.J.\ Buckholtz for making the preprint publicly available.
This work used PDG 2024 constants, Planck 2018 cosmological parameters, and
the MCXC/XMM-Newton galaxy cluster catalogs.
\end{acknowledgments}

% ──────────────────────────────────────────────
\bibliography{refs}
% ──────────────────────────────────────────────

\end{document}
